\documentclass[11pt]{article}
\usepackage[a4paper,margin=27mm]{geometry}
\usepackage{amsmath,amssymb,hyperref}
\hypersetup{colorlinks=true,urlcolor=blue}
\title{Infinitely many constructed Hadamard orders\\A disproof of Conjecture 00000006331}
\author{ziangni-sys}\date{10 October 2026}
\begin{document}\maketitle
\begin{abstract}
The finiteness clause for the Hadamard construction spectrum is false: Sylvester's elementary doubling construction produces Hadamard matrices of order $2^n$ for every $n\ge0$. We prove the sign and row-orthogonality conditions for the actual recursively constructed matrices and show that their orders are unbounded. The accompanying Lean project proves these facts for all $n$, including the infinitude of the attainable-order spectrum.
\end{abstract}

\section{The clause being disproved}
The English statement says that the smallest missing order is given by gap-filling paths of the construction spectrum ``and the spectrum is finite.'' The Chinese statement likewise asserts that the construction spectrum is finite. We use the standard order-spectrum meaning
\[
\mathcal H=\{m\in\mathbb N:\text{a Hadamard matrix of order }m\text{ can be constructed}\}.
\]
In particular, an explicit matrix construction contributes its actual order to this spectrum. The proof below disproves finiteness under this attainable-order interpretation. The source supplies no separate definition restricting the spectrum to a particular finite search range or a finite set of primitive building blocks.

This is not a claim about the set of \emph{missing} orders, and it does not settle existence at every positive order divisible by four. Nor does it classify Paley constructions or the source's unspecified gap-filling paths. One false conjunct suffices to disprove the stated conjunction.

\section{Matrices of every power-of-two order}
A Hadamard matrix of order $m$ is a square matrix $H$ with entries in $\{1,-1\}$ satisfying
\[
\sum_{k=1}^{m}H_{ik}H_{jk}=m\,\mathbf 1_{i=j}.
\]
Define actual integer matrices recursively by
\[
H_0=(1),\qquad
H_{n+1}=\begin{pmatrix}H_n&H_n\\H_n&-H_n\end{pmatrix}.
\]
Inductively, $H_n$ has $2^n$ rows and columns and every entry is $1$ or $-1$.

For completeness, index its rows and columns by $I_0=\{*\}$ and $I_{n+1}=\{0,1\}\times I_n$, and set
\[
H_{n+1}((a,x),(b,z))=(-1)^{ab}H_n(x,z).
\]
The inner product of two rows is
\begin{align*}
&\sum_{b\in\{0,1\}}\sum_{z\in I_n}
 H_{n+1}((a,x),(b,z))H_{n+1}((c,y),(b,z))\\
&\quad=\left(\sum_{b\in\{0,1\}}(-1)^{ab}(-1)^{cb}\right)
 \left(\sum_{z\in I_n}H_n(x,z)H_n(y,z)\right)\\
&\quad=(2\mathbf1_{a=c})(2^n\mathbf1_{x=y})
 =2^{n+1}\mathbf1_{(a,x)=(c,y)}.
\end{align*}
The base case is immediate, so this proves the exact Hadamard identity for every $n$. Equivalently, $H_nH_n^{\mathsf T}=2^n I_{2^n}$.

\section{The spectrum cannot be finite}
Every $2^n$ belongs to $\mathcal H$ by the explicit construction. Given any $N\in\mathbb N$, choose $n=N+1$. The elementary inequality $n<2^n$ gives
\[
N<2^{N+1}\in\mathcal H.
\]
Thus $\mathcal H$ is unbounded, and hence infinite. Even if one discards the orders one and two and restricts to positive multiples of four, the subfamily $2^{n+2}$ is still unbounded. This directly contradicts the finite construction-spectrum clause.

\section{Formalization and reproduction}
The Lean project defines the recursive index types and the integer matrices themselves. It proves their cardinalities, sign-valued entries, and full row-inner-product identities by induction. An explicit finite equivalence reindexes the matrices onto $\operatorname{Fin}(2^n)$, giving precisely the ordinary finite square-matrix definition of Hadamard order. Finally, it proves that every bound is exceeded by an actual Hadamard order, and contradicts the maximum of any proposed finite spectrum.

The principal theorems are \texttt{entries\_sign}, \texttt{row\_inner}, \texttt{sylvester\_is\_hadamard}, and \texttt{construction\_spectrum\_infinite}. No matrix existence or orthogonality hypothesis is assumed. There are no proof placeholders, custom axioms, or native evaluation shortcuts.

Use Lean 4.19.0 and run \texttt{lake build} in the \texttt{lean/} directory. Mathlib is pinned to commit
\begin{center}\texttt{c44e0c8ee63ca166450922a373c7409c5d26b00b}.\end{center}
The final build prints the axiom dependencies of the principal theorems. Only standard Lean/Mathlib logical axioms are permitted. The argument is entirely symbolic; no numerical search or external certificate is required.
\end{document}
